\section{Comparative Analysis}
This section positions the proposed AirGlove prototype against existing commercial air-mouse products and relevant research prototypes. The goal is to clarify what the system already does well, where it currently falls short, and which constraints are most likely to affect real-world usability and adoption. The analysis is structured in three parts: first, a SWOT evaluation summarizes internal strengths and limitations alongside external opportunities and risks; second, a structured comparison table consolidates key differences in form factor, sensing approach, interaction capabilities, connectivity, and practical constraints; finally, a brief research and industry context highlights trends that inform realistic improvement paths for future iterations.

\subsection{SWOT Analysis}

In order to better understand the current state of midair pointing devices and wearable input technologies, existing commercial solutions and research prototypes were analyzed. The market currently includes handheld air mice, presenter remotes, and experimental glove-based input devices, which demonstrate both the practicability and the limitations of midair cursor control. Based on this analysis, a SWOT (Strengths, Weaknesses, Opportunities, Threats) evaluation was conducted to assess the proposed air mouse glove solution using ESP32 and MPU6050/MPU6500 sensors. The full comparative analysis is provided in Appendix~\ref{appendix:comparative_table}, while Table~\ref{tbl:swot} summarizes the main strengths, weaknesses, opportunities, and threats of the proposed design.


\subsection{Comparative Analysis}

As shown in Table \ref{tbl:comparative_air_mouse}, the comparative analysis highlights several important distinctions across the evaluated solutions. In terms of form factor and wearability, the proposed glove shares the wearable category with experimental concepts such as Ion and Bellcos, as well as research prototypes, but differentiates itself from handheld remotes and presenter devices by eliminating the need to grip or hold any physical object.

Regarding interaction method, commercial air mice rely on tilt-based motion but offer no gesture vocabulary beyond cursor movement and button clicks, whereas the proposed design inherits gesture support (scroll, zoom, pause mode) from its research-prototype lineage \cite{cornell_airmouse_project,acm_airmouse_interface} while adding Bluetooth HID connectivity characteristic of commercial devices. With respect to sensor technology, all IMU-based solutions share the same fundamental sensing principle (gyroscope and accelerometer fusion) but differ substantially in the quality: commercial products ship with factory-tuned firmware, while the proposed prototype requires custom filter implementation. This represents both the main current weakness of the prototype and its primary long-term advantage, since the filtering strategy is fully adjustable and can be iteratively improved.

On connectivity, the proposed glove is designed to operate primarily as a native Bluetooth HID device without a USB dongle, which improves portability compared to 2.4 GHz receiver-dependent commercial remotes. A 2.4 GHz variant remains a possible extension for environments where Bluetooth is restricted or unreliable. Regarding precision and stability, the proposed design is rated as tunable, honestly reflecting the current prototype stage; commercial devices rated medium or medium--high are factory-optimised but cannot be modified, whereas the proposed glove's precision ceiling is determined by the sophistication of the algorithms applied rather than by fixed hardware constraints. 

Battery life is a dimension routinely overlooked in product-level comparisons of air mice, yet it is critically important for wearable devices intended for real-world use. Commercial devices suffer a well-documented limitation in this regard: user reviews consistently cite short operational duration as a key frustration \cite{reddit_airmouse_discussion}, and none of the handheld or experimental wearable solutions evaluated offer a transparent power management strategy. The proposed design, powered by a LiPo battery connected to the ESP32, does not resolve this issue at the prototype stage; however, its open architecture allows power consumption to be actively managed through deep-sleep modes and duty-cycling, an optimisation path unavailable in any closed commercial product.

Similarly, customisability is a differentiating factor that commercial solutions cannot offer by design. The proposed glove is unique among the evaluated solutions in being fully open: every layer of the system, from sensor calibration and gesture mapping to Bluetooth HID report descriptors, is accessible and modifiable. This is significant not only for ongoing development but also for the research and accessibility use cases the device targets, where gesture vocabularies and sensitivity profiles may need to be adapted per user or per deployment context.

Overall, the picture that emerges is nuanced. The proposed design does not yet match commercial devices in terms of polish, battery life, or out-of-box precision. However, it outperforms them in portability, wearability, customisability, and the potential ceiling of its pointing performance, advantages being only realisable through continued algorithmic development and iterative hardware refinement.

\subsection{Industry Trends and Research Context}

The proposed air mouse glove is positioned within a specific and well-defined niche: IMU-based wearable devices used as pointing and cursor-control interfaces. This section surveys relevant research and commercial developments within that niche, rather than the broader gesture recognition industry, which includes unrelated application domains such as automotive infotainment, touchscreen panels, and computer-vision systems.


\subsubsection{IMU Wearables as Pointing Devices: Research Trends}

Research interest in IMU-based wearable pointing devices has grown steadily within the HCI community. A representative recent example is \textit{MouseRing}, presented at ACM CHI~2024 \cite{acm_airmouse_interface}, which proposes an IMU ring-shaped device enabling continuous finger-based cursor control on unmodified surfaces, targeting always-available interaction in AR/VR and large-screen environments. The work explicitly identifies the same core design tension faced by the proposed glove: the need to retain the comfort and precision of conventional pointing while eliminating the physical and environmental constraints of surface-dependent devices.

Similarly, research into IMU-based assistive mouse controllers, such as the \textit{Auxilio} system, which uses a low-cost commercial IMU to enable cursor control for users with upper-limb disabilities, demonstrates that the principle of mapping IMU orientation to cursor movement is already validated in peer-reviewed HCI literature \cite{arxiv_gesture_recognition}. These works confirm that the sensing approach underlying the proposed air mouse glove is technically sound and aligned with active research directions.

\subsubsection{IMU as the Preferred Sensor Modality for Wearable Pointing}

Within the wearable pointing-device domain, IMUs have emerged as the sensor of choice due to their low power consumption, compact form factor, low cost, and independence from line-of-sight or lighting constraints, limitations that affect camera-based alternatives. Industry sources confirm that IMUs are increasingly embedded in wrist-worn and finger-worn devices for gesture input, and that their effectiveness is directly tied to the quality of the algorithms processing their output \cite{gesture_market}.

Critically, research consistently demonstrates that hardware selection alone is insufficient to determine pointing quality; the choice of sensor fusion strategy, whether Madgwick, Mahony, Complementary, or Kalman filtering, has a substantial effect on stability and drift. A 2024 study published in the \textit{International Journal of Reconfigurable and Embedded Systems} confirmed that IMU gesture classifiers using full orientation representations achieve accuracy above 99\% under embedded resource constraints, while reduced feature sets maintain accuracy above 90\% \cite{gesture_hci_research}. These findings directly support the algorithmic approach adopted in the proposed design, in which sensor fusion and calibration routines are treated as core engineering priorities rather than secondary concerns.

\subsubsection{Accessibility and the Case for Surface-Independent Input}

A recurring motivation in the IMU wearable pointing literature is accessibility: enabling computer interaction for users who cannot comfortably use a conventional mouse or touchpad. The \textit{Auxilio} project specifically targets users with upper-limb disabilities by replacing the desktop mouse with an IMU-driven interface, demonstrating that cursor mapping from body-worn inertial sensors is a clinically viable approach.

More broadly, the HCI research community has documented demand for input solutions that are portable, surface-independent, and ergonomically undemanding, which are requirements that align directly with the intended use cases of the proposed glove, namely presentation control and accessibility. The Wearable Devices and Doublepoint collaboration further illustrates that this design space is attracting commercial investment: their \textit{WowMouse} platform converts an Apple Watch into an air gesture pointing device using on-wrist IMU data and embedded algorithms, demonstrating commercial viability for the same underlying concept.

\subsubsection{Future Directions: Miniaturisation and On-Device AI}

Two related trends are likely to influence the next generation of IMU-based pointing devices. The first is continued miniaturisation of IMU hardware. New “smart IMU” solutions, such as the Bosch BHI360, can run gesture-related algorithms directly on the sensor, which supports smaller designs and lower power consumption.

For the proposed glove, this means that the current ESP32 + MPU6050/6500 setup, while suitable for prototyping, could later be replaced with a more integrated module that offers better drift behaviour and reduced energy use. The second trend is progress in TinyML, where machine-learning models are designed to run efficiently on microcontrollers. Recent work shows that on-device gesture classification can fit within the memory and compute limits of platforms similar to the ESP32 \cite{arxiv_airmouse_future}. This suggests a possible future upgrade path rather than a current project goal: the present system will rely on rule-based calibration, while a later iteration could add lightweight trained classifiers to improve robustness across different users and reduce sensitivity to unintended movements, all while keeping processing on-device.